\begin{lemma}[From the gradient's Campanato decay on small balls to $C^{1+\beta}$]
Let $n\ge2$, $p>1$, $0<\lambda_0<p$, $c_{20}>0$, $0<\rho_0\le1$ and $W\ge0$.
Then there is a constant $N\ge0$, depending only on
$n,p,\lambda_0,c_{20},\rho_0$ and $W$, with the following property.

Let $\Omega,\Omega'\subseteq\R^n$ be open with
$\overline{\Omega'}\subseteq\Omega$, let $\Omega'$ be a bounded Lipschitz domain
(\noderef{IsBoundedLipschitzDomain}), and assume the cited Campanato embedding
(C5) at the exponent $p$ on $\Omega'$ (\noderef{CitedCampanatoEmbedding}).  Let
$u\in W^{1,p}(\Omega)$ with weak gradient $g=\nabla u$ (\noderef{MemW1pOn})
satisfy $\int_\Omega|\nabla u|^p\,dx\le W^p$.  Suppose that every ball of
radius at most $\rho_0$ centred in $\Omega'$ lies in $\Omega$,
\[
B_\rho(x_0)\subseteq\Omega
\qquad\bigl(x_0\in\Omega',\ 0<\rho\le\rho_0\bigr),
\tag{C}
\]
and suppose the decay
\[
\int_{B_\rho(x_0)}\bigl|\nabla u-(\nabla u)_{B_\rho(x_0)}\bigr|^p\,dx
\le c_{20}\,\rho^{\,n+\lambda_0}
\qquad\bigl(x_0\in\Omega',\ 0<\rho\le\rho_0\bigr),
\tag{D}
\]
where $(\nabla u)_{B}=|B|^{-1}\int_B\nabla u$ is the mean of the gradient over
$B$.

\medskip
\noindent\textbf{On the containment $(C)$.}  The radius $\rho_0$ is bound in the
head, before $\Omega$ and $\Omega'$, so no relation between $\rho_0$ and
$d(\Omega',\partial\Omega)$ is available from the other hypotheses; $(C)$ is
that relation, stated at the only strength the proof uses.  It is also exactly
what the paper guarantees at this point: Lemma~4.1 produces its radius $R_1$
``depending only on $n,p,\lambda,\Lambda,a,\nu,M$ and
$d(\Omega',\partial\Omega)$'' and then integrates over the \emph{full} ball
$B_R$ for $R\le R_1$ and $x_0\in\Omega'$ (thesis l.~1162--1168), and making the
radius depend on the separation is how it keeps those balls inside $\Omega$.
The hypothesis is needed because $(D)$ and the first display of Step 4 integrate
over the whole ball $B_\rho(x_0)$, while $u$ and $g$ are controlled only on
$\Omega$: without $(C)$ the integrand of $(D)$ need not be integrable on
$B_\rho(x_0)$ at all.  It is consumed in exactly one place, the enlargement of
the domain of integration in Step 4.  Both consumers supply it at once from
their own separation hypothesis, the radius they pass as $\rho_0$ being at most
half the radius for which they have already recorded
$\overline B_R(x_0)\subseteq\Omega$.  Then
\begin{itemize}
\item[(i)] $\|D_iu\|_{\mathcal L^{p,\,n+\lambda_0}(\Omega')}\le N$ for every
  $i=1,\dots,n$ (\noderef{CampanatoBoundOn}); and
\item[(ii)] there are $w\colon\R^n\to\R$ and $g_w\colon\R^n\to\R^n$ with $u=w$
  almost everywhere on $\Omega'$, $g=g_w$ almost everywhere on $\Omega'$, and
  $\mathrm{MemC1HolderOn}\;n\;\beta\;\Omega'\;\overline{\Omega'}\;w\;g_w$ at
  $\beta=\lambda_0/p$ (\noderef{MemC1HolderOn} at the pair of sets
  $(\Omega',\overline{\Omega'})$) --- that is, $w$ is differentiable at every
  point of $\Omega'$ with $D_iw=(g_w)_i$ there for every $i$, $g_w$ is
  continuous on $\overline{\Omega'}$, and
  $\|g_w(x)-g_w(y)\|\le C\operatorname{dist}(x,y)^{\beta}$ for all
  $x,y\in\overline{\Omega'}$ and some $C\ge0$.
\end{itemize}
\end{lemma}
\begin{proof}
\emph{Step 1: the constant, chosen before the data.}  Put
\[
N=2\,\max\Bigl\{c_{20},\ W^p\rho_0^{-(n+\lambda_0)}\Bigr\}^{1/p}\ \ge0 ,
\]
a function of $n,p,\lambda_0,c_{20},\rho_0$ and $W$ alone; it is nonnegative
because both entries of the maximum are nonnegative ($c_{20}>0$, and
$W^p\ge0$, $\rho_0^{-(n+\lambda_0)}>0$).  Raising to the $p$-th power,
\[
N^p=2^{\,p}\max\Bigl\{c_{20},\ W^p\rho_0^{-(n+\lambda_0)}\Bigr\}
\ \ge\ 2^{\,p}c_{20},
\qquad
N^p\ \ge\ 2^{\,p}W^p\rho_0^{-(n+\lambda_0)} .
\tag{1}
\]
Two further consequences, recorded now.  First, $n+\lambda_0>0$, since $n\ge2$
and $\lambda_0>0$; so $0<\rho_0\le1$ gives $\rho_0^{-(n+\lambda_0)}\ge1$ and
therefore, by the second half of $(1)$,
\[
N^p\ \ge\ 2^{\,p}W^p,\qquad\text{hence}\qquad N\ \ge\ 2W\ \ge\ W
\tag{2}
\]
(taking $p$-th roots, both sides being nonnegative, and using $W\ge0$).
Second, $\beta=\lambda_0/p$ satisfies $0<\beta<1$, because $0<\lambda_0<p$.

Fix now $\Omega,\Omega',u,g$ satisfying all the hypotheses.  Note
$\Omega'\subseteq\overline{\Omega'}\subseteq\Omega$.

\emph{Step 2: the components of the gradient lie in $L^p(\Omega')$.}  By
\noderef{MemW1pOn} the field $g$ lies in $L^p(\Omega)$; since $\Omega'$ is open,
hence measurable, and $\Omega'\subseteq\Omega$, it lies in $L^p(\Omega')$ as
well, $L^q$ membership being inherited by measurable subsets (Preamble,
H\"older's inequality on a set of finite measure, the first of the two transport
facts recorded there).  Fix $i$.  The function $x\mapsto g(x)_i$ is the
composition of $g$ with the $i$-th coordinate projection of $\R^n$, which is
linear and continuous, so it is almost everywhere strongly measurable on
$\Omega'$; and $|g(x)_i|\le|g(x)|$ pointwise, the modulus of one coordinate of a
vector of $\R^n$ being at most its Euclidean norm.  Hence
$|g(\cdot)_i|^p\le|g|^p$ pointwise, and $|g|^p$ is integrable on $\Omega'$
because $g\in L^p(\Omega')$; so $|g(\cdot)_i|^p$ is integrable on $\Omega'$ by
the Preamble's domination clause (a), i.e.
\[
D_iu=g(\cdot)_i\in L^p(\Omega')
\tag{3}
\]
--- membership in $L^p$ of a set, for the finite exponent $p$, being almost
everywhere strong measurability together with integrability of the $p$-th power
of the modulus.  The same argument gives $D_iu\in L^p(S)$ for every measurable
$S\subseteq\Omega'$.

\emph{Step 3: the mean-replacement bound.}  Let $S\subseteq\R^n$ be measurable
with $0<|S|<\infty$, let $h\in L^p(S)$ be \emph{real}-valued, let $c\in\R$, and
write $h_S=|S|^{-1}\int_Sh$.  Then
\[
\int_S|h-h_S|^p\,dx\le2^{\,p}\int_S|h-c|^p\,dx .
\tag{M}
\]
Indeed, replacing $h$ by $h-c$ --- which replaces $h_S$ by $h_S-c$, by
additivity of the integral and $|S|^{-1}\int_Sc=c$, and leaves $h-h_S$
unchanged --- it is enough to prove this for $c=0$.  By convexity of
$t\mapsto t^p$ on $[0,\infty)$, $|h-h_S|^p\le2^{p-1}(|h|^p+|h_S|^p)$ pointwise,
so it suffices to bound $|h_S|$.

\emph{For that, $\bigl|\int_Sh\bigr|\le\int_S|h|$, and it is obtained by
splitting $h$ into its nonnegative parts rather than from a norm inequality for
the Bochner integral.}  The Preamble's norm-inequality clause is stated in
\emph{interval}-integral form and does not apply to an integral over
$S\subseteq\R^n$; what the Preamble does declare over a general measure space is
the domination clause and H\"older's inequality for nonnegative functions, and
those suffice.  First, $h\in L^1(S)$: by the H\"older item's own consequence,
$L^{p}(S)\subseteq L^1(S)$ because $|S|<\infty$ and $p\ge1$.  Next put
$h^+=\max\{h,0\}$ and $h^-=\max\{-h,0\}$.  Both are measurable with
$0\le h^{\pm}\le|h|$, so both are integrable on $S$ by the Preamble's
domination clause (a), which is stated for an arbitrary measure space.  At every
point $h=h^+-h^-$ and $|h|=h^++h^-$ --- check the two cases $h\ge0$ and $h<0$
--- so by additivity of the integral
\[
\int_Sh=\int_Sh^+-\int_Sh^- ,
\qquad
\int_S|h|=\int_Sh^++\int_Sh^- .
\]
Both $\int_Sh^+$ and $\int_Sh^-$ are nonnegative, the integral of a nonnegative
integrable function over a general measure space being nonnegative --- the fact
the Preamble's H\"older item already presupposes, since that item is stated for
$f,g\ge0$ on an arbitrary measure space.  For nonnegative reals $P,Q$ one has
$|P-Q|\le P+Q$, which gives $\bigl|\int_Sh\bigr|\le\int_S|h|$ exactly.

Combining with H\"older's inequality on the finite-measure set $S$ (Preamble,
the consequence $\int_S|h|\le|S|^{1-1/p}(\int_S|h|^p)^{1/p}$),
\[
|h_S|=\Bigl|\frac1{|S|}\int_Sh\Bigr|\le\frac1{|S|}\int_S|h|
\le\Bigl(\frac1{|S|}\int_S|h|^p\Bigr)^{1/p} .
\]
Hence $|S|\,|h_S|^p\le\int_S|h|^p$, and integrating the pointwise bound over $S$
gives
$\int_S|h-h_S|^p\le2^{p-1}\bigl(\int_S|h|^p+\int_S|h|^p\bigr)=2^p\int_S|h|^p$,
which is (M) at $c=0$.

\emph{Step 4: part \textup{(i)}.}  Fix $i$, let $x_0\in\Omega'$ and
$0<\rho<\operatorname{diam}\Omega'$, and put $S=\Omega'\cap B_\rho(x_0)$, a
measurable set with $|S|\le|B_\rho(x_0)|<\infty$.  If $|S|=0$ then the integral
on the left of the Campanato inequality vanishes, being an integral over a null
set, while its right-hand side $N^p\rho^{\,n+\lambda_0}$ is nonnegative, and
there is nothing to prove.  So assume $|S|>0$; then (M) applies on $S$ with
$h=D_iu$, which lies in $L^p(S)$ by Step 2.

If $\rho\le\rho_0$, apply (M) with
$c=\bigl((\nabla u)_{B_\rho(x_0)}\bigr)_i$, and then enlarge the domain of
integration and pass from the component to the vector:
\[
\int_S\bigl|D_iu-(D_iu)_S\bigr|^p
\le2^{\,p}\!\int_S\bigl|D_iu-c\bigr|^p
\le2^{\,p}\!\int_{B_\rho(x_0)}\bigl|\nabla u-(\nabla u)_{B_\rho(x_0)}\bigr|^p
\le2^{\,p}c_{20}\,\rho^{\,n+\lambda_0}
\le N^p\rho^{\,n+\lambda_0},
\]
the second step because $S\subseteq B_\rho(x_0)$ and the integrand is
nonnegative with
$|D_iu-c|=\bigl|\bigl(\nabla u-(\nabla u)_{B_\rho(x_0)}\bigr)_i\bigr|
\le\bigl|\nabla u-(\nabla u)_{B_\rho(x_0)}\bigr|$ pointwise, the third by (D),
and the last by the first half of $(1)$.

\emph{The second step is where $(C)$ is consumed, and it is the only place.}
Enlarging the domain of integration from $S$ to $B_\rho(x_0)$ for a nonnegative
integrand requires that integrand to be integrable on the larger set, and
$\nabla u$ is controlled only on $\Omega$.  By $(C)$, applied at the present
$x_0\in\Omega'$ and the present $\rho\le\rho_0$, we have
$B_\rho(x_0)\subseteq\Omega$; so $\nabla u\in L^p(B_\rho(x_0))$ by
\noderef{MemW1pOn} and the inheritance of $L^p$ membership by measurable
subsets (Preamble, H\"older's inequality on a set of finite measure, the first
of the two transport facts recorded there), hence
$\nabla u-(\nabla u)_{B_\rho(x_0)}\in L^p(B_\rho(x_0))$ --- the subtracted
vector being constant and $B_\rho(x_0)$ of finite measure --- and therefore
$\bigl|\nabla u-(\nabla u)_{B_\rho(x_0)}\bigr|^p$ is integrable on
$B_\rho(x_0)$, which is what the enlargement needs.

If instead $\rho>\rho_0$, apply (M) with $c=0$:
\[
\int_S\bigl|D_iu-(D_iu)_S\bigr|^p\le2^{\,p}\!\int_S|D_iu|^p
\le2^{\,p}\!\int_\Omega|\nabla u|^p\le2^{\,p}W^p
\le2^{\,p}W^p\rho_0^{-(n+\lambda_0)}\rho^{\,n+\lambda_0}
\le N^p\rho^{\,n+\lambda_0},
\]
the second step because $S\subseteq\Omega$ and $|D_iu|^p\le|\nabla u|^p$
pointwise (Step 2) with $|\nabla u|^p$ integrable on $\Omega$, the third by the
hypothesis $\int_\Omega|\nabla u|^p\le W^p$, the fourth because
$\rho_0^{-(n+\lambda_0)}\rho^{\,n+\lambda_0}=(\rho/\rho_0)^{\,n+\lambda_0}\ge1$
as $\rho>\rho_0>0$ and $n+\lambda_0>0$, and the last by the second half of
$(1)$.

So in both regimes
$\int_S|D_iu-(D_iu)_S|^p\le N^p\rho^{\,n+\lambda_0}$, which together with
$(3)$ is exactly
$\|D_iu\|_{\mathcal L^{p,\,n+\lambda_0}(\Omega')}\le N$
(\noderef{CampanatoBoundOn}), the two clauses of that predicate being the
$L^p(\Omega')$ membership and the displayed bound.  That is (i).

\emph{Step 5: the \textup{(C5)} step.}  Apply the assumed cited Campanato
embedding (C5) at the exponent $p$ on $\Omega'$
(\noderef{CitedCampanatoEmbedding}) with $\lambda=n+\lambda_0$ and with the
\emph{same} $N$ as in (i).  Its Lipschitz antecedent is the hypothesis that
$\Omega'$ is a bounded Lipschitz domain (\noderef{IsBoundedLipschitzDomain}).
Read against that node's statement, it binds three guards and then \emph{two}
hypotheses on the function, and both must be supplied at that one $N$:
\begin{itemize}
\item the guards $n<\lambda$ and $\lambda<n+p$, which hold because
  $0<\lambda_0<p$, and $0\le N$, which Step 1 records;
\item $\mathcal L^{p,n+\lambda_0}(\Omega')$ with norm at most $N$
  (\noderef{CampanatoBoundOn}) for the function in question, which for $D_iu$ is
  exactly (i); and
\item \emph{the Morrey hypothesis at the exponent $0$}, i.e.\
  $L^{p,0}(\Omega')$ with norm at most the \emph{same} $N$
  (\noderef{MorreyBoundOn}).  Unfolding \noderef{MorreyBoundOn} at $\lambda=0$,
  this asks for $D_iu\in L^p(\Omega')$, which is $(3)$, together with
  \[
  \int_{\Omega'\cap B_\rho(x)}|D_iu|^p\,dx\le N^p\rho^{\,0}=N^p
  \qquad\bigl(x\in\Omega',\ 0<\rho<\operatorname{diam}\Omega'\bigr),
  \]
  using $\rho^{\,0}=1$ for $\rho>0$.  This holds because
  $\int_{\Omega'\cap B_\rho(x)}|D_iu|^p\le\int_\Omega|\nabla u|^p\le W^p\le N^p$
  --- the first step as in Step 4, the second by hypothesis, and the third by
  $(2)$ together with $W\ge0$ and the monotonicity of $t\mapsto t^p$ on
  $[0,\infty)$.  \emph{This is not the same as ``$D_iu\in L^p(\Omega')$'': the
  node forces one constant to serve both the Campanato and the Morrey bound, so
  the inequality $W\le N$ of $(2)$ is load bearing and is why $N$ was defined
  with the factor $2$.}
\end{itemize}
The guards and the two hypotheses do not depend on $i$, and that node binds its
constant before the function, so one constant $N_h\ge0$ serves all $i$.  For
each $i$ it produces a function $w^{(i)}$ with $D_iu=w^{(i)}$ almost everywhere
on $\Omega'$, with $w^{(i)}\in C^{\beta}(\overline{\Omega'})$
(\noderef{MemHolderOn}) at the exponent $(\lambda-n)/p=\lambda_0/p=\beta$, and
with
\[
|w^{(i)}(x)-w^{(i)}(y)|\le N_h\operatorname{dist}(x,y)^{\beta}
\qquad\bigl(x,y\in\overline{\Omega'}\bigr).
\]

Put $g_w(x)=\bigl(w^{(1)}(x),\dots,w^{(n)}(x)\bigr)$.  Then:
\begin{itemize}
\item $g=g_w$ almost everywhere on $\Omega'$.  For each of the finitely many $i$
  the set where $D_iu\ne w^{(i)}$ is null, so the union of these $n$ sets is
  null, and off that union all $n$ coordinates of $g$ and of $g_w$ agree, hence
  the vectors agree.
\item $g_w$ is continuous on $\overline{\Omega'}$: each $w^{(i)}$ is continuous
  there, by the first clause of its \noderef{MemHolderOn} assertion, and a map
  into $\R^n$ whose $n$ coordinate functions are continuous on a set is
  continuous on that set.
\item $\|g_w(x)-g_w(y)\|\le\sqrt n\,N_h\operatorname{dist}(x,y)^{\beta}$ for
  $x,y\in\overline{\Omega'}$: the Euclidean norm of a vector of $\R^n$ is at
  most $\sqrt n$ times the largest modulus of its coordinates, and each
  coordinate of $g_w(x)-g_w(y)$ is $w^{(i)}(x)-w^{(i)}(y)$, of modulus at most
  $N_h\operatorname{dist}(x,y)^{\beta}$.
\end{itemize}
These are the second and third clauses of
$\mathrm{MemC1HolderOn}\;n\;\beta\;\Omega'\;\overline{\Omega'}\;w\;g_w$
(\noderef{MemC1HolderOn}), with the constant $\sqrt n\,N_h\ge0$, for whatever
$w$ the next step produces.

\emph{Step 6: the first clause, i.e.\ the differentiable representative.}  What
remains is a single $w$ with $u=w$ almost everywhere on $\Omega'$ which is
differentiable at every point of $\Omega'$ with $\nabla w=g_w$ there.

\emph{First, $u\in W^{1,p}_{\mathrm{loc}}(\Omega')$.}  Apply
\noderef{WeakGradientRestrict} twice; its four conclusions are conditionals
under one binder block, so reaching $W^{1,p}_{\mathrm{loc}}$ on $\Omega'$ takes
two instantiations.
\begin{itemize}
\item At $n$, $q=p$, the sets $(\Omega,\Omega')$ and the pair $(u,g)$.  Its
  binders are $1\le n$, from $2\le n$; $1\le p$, from $1<p$;
  $\mathrm{IsOpen}\,\Omega$ and $\mathrm{IsOpen}\,\Omega'$, both hypotheses
  here; and $\Omega'\subseteq\Omega$, recorded at the end of Step 1.  Its
  \emph{first} conclusion turns the hypothesis $u\in W^{1,p}(\Omega)$ with weak
  gradient $g$ (\noderef{MemW1pOn}) into $u\in W^{1,p}(\Omega')$ with the same
  weak gradient.
\item Again at $n$, $q=p$, now the sets $(\Omega',\Omega')$ and the same pair.
  Its binders are $1\le n$ and $1\le p$ as before, $\mathrm{IsOpen}\,\Omega'$
  twice, and $\Omega'\subseteq\Omega'$, which is trivial.  Its \emph{third}
  conclusion turns the membership just obtained into
  $u\in W^{1,p}_{\mathrm{loc}}(\Omega')$ with weak gradient $g$
  (\noderef{MemW1pLocOn}).
\end{itemize}

\emph{Then the representative.}  Apply \noderef{WeakGradientContinuousRep} with
the present $n$, with $q=p$, with the open set $\Omega'$, with $u$, with $g$ and
with $g_c=g_w$.  Its hypotheses, each discharged by name:
\begin{itemize}
\item $1\le n$, from $2\le n$;
\item $1\le p$, from $1<p$;
\item $\mathrm{IsOpen}\,\Omega'$, a hypothesis here;
\item $u\in W^{1,p}_{\mathrm{loc}}(\Omega')$ with weak gradient $g$
  (\noderef{MemW1pLocOn}), which is the second bullet above;
\item $g_w$ continuous on $\Omega'$: Step 5 gives $g_w$ continuous on
  $\overline{\Omega'}$, and continuity on a set restricts to any subset, with
  $\Omega'\subseteq\overline{\Omega'}$;
\item $g=g_w$ almost everywhere on $\Omega'$, which is Step 5.
\end{itemize}
It produces a function $w$ with $u=w$ almost everywhere on $\Omega'$ and with
$w$ differentiable at every point $x\in\Omega'$ and $D_iw(x)=g_w(x)_i$ for every
$i$.  That is exactly the \emph{first} clause of
$\mathrm{MemC1HolderOn}\;n\;\beta\;\Omega'\;\overline{\Omega'}\;w\;g_w$
(\noderef{MemC1HolderOn}), read at the pair of sets
$(\Omega',\overline{\Omega'})$, whose first clause ranges over the first set.

\emph{Assembly.}  Conjoining that first clause with the second and third
clauses established in Step 5 gives
$\mathrm{MemC1HolderOn}\;n\;\beta\;\Omega'\;\overline{\Omega'}\;w\;g_w$ with
$\beta=\lambda_0/p$; together with $u=w$ almost everywhere on $\Omega'$ and
$g=g_w$ almost everywhere on $\Omega'$, this is (ii).  The constant $N$ was
fixed in Step 1, before $\Omega,\Omega',u$ and $g$, which is what the statement
asserts of it.
\end{proof}
